\subsection{The Original Express Clocks}
\label{sec:original-express}

The breakpoint could wait for an inspector's reply. A running clock poses
a different problem: how can a new request affect a handler that has
not finished? An ordinary object finishes one script before starting the next queued
request. This creates a problem for a clock whose \code{:start} script
loops until another request sets \code{stop}: an ordinary stop request
waits behind the very script it must stop. Express messages provide an
alternative dispatch policy, not faster transport or operating-system
threads. Read \filepath{examples/misc/express-ex.abclr2} with
\filepath{lib/express.abclr2}; its older companion uses
\filepath{lib/express-old.abclr2}.

\paragraph{Configuration and the application first.}
The original group installs both halves of the protocol:

\wholelisting[firstline=9,lastline=12]{examples/misc/express-ex.abclr2}

The evaluator recognizes express sends; the metaobject generator equips
receivers to accept them. Both \code{clock} and \code{test} belong to this
group. A keyword in a send does not by itself change the primary evaluator
or the receiver's ordinary queue policy. The top-level only sends an
ordinary request to \code{test}; express sends occur inside that object.

The clock starts with \code{time = 0} and \code{stop = nil}.
Its \code{:start} script resets only \code{stop}, not \code{time}:

\wholelisting[firstline=19,lastline=31]{examples/misc/express-ex.abclr2}

In Common Lisp, \code{(do () (stop) ...)} has no loop variables and tests
\code{stop} before each iteration. \code{(dotimes (i 5))} runs five empty
iterations; it is neither a delay in seconds nor a clock measurement.
Each completed pair adds two. The intermediate value between increments
is odd, so excluding express observation of that intermediate state is
the point of \code{atomic}. Without it, an odd sample is permitted, not
guaranteed: a particular schedule may still produce only even samples.

\wholelisting[firstline=44,lastline=50]{examples/misc/express-ex.abclr2}

The comments calling these scripts ``express'' do not declare a special
kind of script. The \emph{send} chooses the mode; an ordinary \code{:peek}
uses the same matching script. \code{:peek} replies with the current count.
\code{:stop} changes a binding but does not reply. Stopping is completed
when the resumed loop notices that binding and replies \code{:finished}.

\par\noindent\begin{minipage}{\linewidth}
\paragraph{Following the controller's two phases.}
The controller supplies the orchestration rather than the clock itself:
\par\vspace*{\medskipamount}

\wholelisting[firstline=55,lastline=71]{examples/misc/express-ex.abclr2}
\end{minipage}\par

Trace the original \code{[test <== :start]} as follows:

\begin{enumerate}
  \item \code{R} is the caller's reply destination. The controller sends
  \code{clock} an ordinary asynchronous \code{:start}, attaching a
  continuation instead of waiting for the loop to finish.
  \item Twenty now-type express \code{:peek} requests each wait for one
  count and print it. They are sequential observations by this controller,
  not twenty simultaneous requests.
  \item A past-type express \code{:stop} requests the state change without
  waiting for a reply that the stop script never sends.
  \item The clock's \code{:finished} reply invokes the continuation.
  Its \code{\_} discards that value and sends \code{Me} a request for the
  second phase, explicitly retaining \code{R}.
  \item The second phase repeats the experiment without the paired
  increments' atomic wrapper. The count continues from the first phase.
  Its completion continuation sends \code{:end} to the controller.
  \item Only \code{:end}'s \code{!:finished} answers the original caller.
\end{enumerate}

Thus the immediate controller script finishing, an express peek replying,
and the entire two-phase test completing are three different events.
Replacing the asynchronous start by a now-type start would prevent this
controller from issuing its peeks and stop. Replacing stop by a now-type
stop would wait for a nonexistent application reply.

\paragraph{How the newer library interrupts evaluation.}
The generator creates a dispatcher, a metaobject, and an evaluator front
end for each receiver. The dispatcher finds the script and constructs
its binding environment over the object's persistent state. The front
end forwards ordinary evaluation to the group's evaluator, but can
retain the next \code{[:do . Context]} request instead of executing it.
This is cooperative interruption at an evaluator boundary: it cannot
stop a host Lisp operation halfway through its machine instructions.

\begin{longtable}{L{0.18\textwidth}L{0.72\textwidth}}
\textbf{Mode} & \textbf{Intended handling of an arriving express request}\\\hline
\endhead
\code{:dormant} & Dispatch it directly; there is no normal script to suspend.\\
\code{:active} & Ask the front end to \code{:grab}; save the normal
evaluation request and its completion destination, then dispatch the express script.\\
\code{:express} & Queue it; one express script does not preempt another.\\
\code{:lock} & Queue it until the atomic region releases the lock.\\
\end{longtable}

Normal arrivals during active, express, or locked execution enter the
normal queue. At an express script's completion, the intended order is
pending express requests, resumption of the saved normal context, and
then queued normal requests. The saved context includes the destination
\code{R} of the evaluator request, not necessarily the application
caller's reply destination. Resuming it must preserve evaluation's
completion path as well as expression and environment.

The \code{atomic} macro in \filepath{src/user.lsp} asks \code{[meta Me]}
for a now-type \code{:lock}, evaluates its body, sends a past-type
\code{:release}, and returns the body's last value. In the newer library,
locking saves the previous mode and sets \code{:lock}; releasing restores
that mode and can start an express request deferred during normal
execution. This is object-local protection against this express protocol,
not a transaction over all objects. The macro uses no unwind-protection:
do not infer automatic release after an error or a non-local exit.

\paragraph{Why the older peek locks too.}
The July version has no grabbing front end or suspended normal context.
It dispatches express scripts without putting them into the normal
message queue, while normal evaluation can continue. Its
\code{express-count} tracks outstanding express executions. A separate
lock object stores waiting reply destinations and grants them in order.
Both a writer and an observer must acquire that same receiver-local lock:

\wholelisting[firstline=41,lastline=46]{examples/misc/express-old-ex.abclr2}

An unlocked read in that version could run between the two protected
increments, because the reader does not participate in the lock protocol.
The newer example deliberately removes \code{atomic} from \code{:peek}:
under its intended suspension policy, the normal writer cannot advance
while the express reader executes, and a reader arriving during a locked
pair must wait. These are different algorithms, not interchangeable files.

Both libraries define the same public generator names and provide the
feature \code{express}. Compare them in separate fresh sessions, explicitly
loading the matching library first. The old example's \code{require}
otherwise names \code{express}, not \code{express-old}; its commented
sample also still names \code{express-ex}. Use the actual old filename:

\begin{codeblock}
\begin{lstlisting}
; New version, in a fresh session:
(inc-lib "lib/express" :use-source t)
(inc "examples/misc/express-ex" :use-source t)

; Old version, in a different fresh session:
(inc-lib "lib/express-old" :use-source t)
(inc "examples/misc/express-old-ex" :use-source t)
\end{lstlisting}
\end{codeblock}

\paragraph{Connect the reading to the first experiment.}
The idle-clock exchange introduced earlier isolates sender configuration,
reply forwarding, and the receiver's express dispatch.
Read \code{:start-non-atomic} next: it starts normal clock work, samples
the value, sends stop, and waits for the clock's completion reply.
Follow the roles of the two reply destinations separately before tracing
the atomic phase.

The queue-handling branches dequeue the selected envelope and bind
\code{[Message Reply Sender]} before dispatch. Compare this with the
saved execution context: an envelope describes an application request,
whereas the saved context describes suspended expression evaluation.
The \code{express-check} sessions provide short examples of both
library versions, queue order, and reply routing.

\paragraph{Check your understanding.}
Why is \code{:stop} sent asynchronously? Why is \code{time} not reset
between phases? Could the non-atomic phase print only even values? Why
does configuring only the clock fail to establish express sending?
Answers: stop has no reply; only \code{stop} is reset; scheduling need
not expose an odd intermediate state; and the sender's evaluator must
translate the extra send argument into the express meta-level protocol.
